\documentclass[10pt,twocolumn]{article}
\usepackage[margin=0.75in,columnsep=0.28in]{geometry}
\usepackage{times}
\usepackage{amsmath}
\usepackage{graphicx}
\usepackage{booktabs}
\usepackage{array}
\usepackage{tikz}
\usepackage{titlesec}
\usepackage{enumitem}
\usepackage[colorlinks=true,linkcolor=blue,citecolor=blue,urlcolor=blue]{hyperref}
\usetikzlibrary{arrows.meta,positioning}

\setlength{\parindent}{1em}
\setlength{\parskip}{0pt}

\titleformat{\section}{\normalfont\large\bfseries}{\Roman{section}.}{0.6em}{}
\titleformat{\subsection}{\normalfont\bfseries}{\Alph{subsection}.}{0.5em}{}
\titlespacing*{\section}{0pt}{8pt}{4pt}
\titlespacing*{\subsection}{0pt}{6pt}{3pt}

\begin{document}

\twocolumn[
\begin{center}
{\LARGE \bfseries Court Vision: Motion Planning for an \\ Obstacle-Aware Ball-Dribbling Mobile Robot \\[4pt]}
{\large RBE 550 Motion Planning --- Project Status Update \\[4pt]}
{\normalsize Rudy Arsenault, Tracy Yang \\
Worcester Polytechnic Institute \\[6pt]}
\end{center}
\vspace{4pt}
\hrule
\vspace{10pt}
]

\section{Progress Overview}
At the Week~8 checkpoint, the project is on schedule against the plan
submitted in the proposal. Objectives~1--3 are implemented and have been
exercised in simulation; Objective~4 (redirect-the-bounce planning) and
Objective~5 (evaluation) remain, as originally scheduled for Weeks~9--13.
Git history shows steady, checkpoint-aligned progress: an initial working
dribble by Week~5 (09-15 -- 09-21), a mobile base that dribbles while
turning by Week~6--7 (09-23), and a working static obstacle course with a
collision-free base planner by the same week, followed by a planner
re-architecture and code-quality pass this week (10-02) ahead of this
submission.

One scope note from the proposal: the division of labor described there
(Yang on dribble dynamics/contact control, Arsenault on obstacle mapping/
base planning) ended up reversed in practice -- Arsenault built and
debugged the first working dribble controller, and Yang built the obstacle
course and the two-stage planner. The underlying technical split (dribble
dynamics vs.\ planning) held; only who owns which half changed. See the
Appendix for each member's specific contributions.

\section{System Implemented To Date}

\subsection{Dribble Dynamics and Contact Model (Objective 1)}
The ball is modeled in PyBullet as a sphere with a floor coefficient of
restitution; between contacts it is a free projectile. At each paddle
contact, outgoing vertical velocity is set by a mirror-law-style rule
\cite{buehler1994}: the commanded paddle velocity is solved so that the
ball's outgoing speed is always $-v_{\text{push}}$ regardless of its
incoming speed, i.e., the hand's effective restitution term cancels out of
the result. A closed-form ballistic model (\texttt{BallisticModel})
predicts flight time analytically from the push speed, floor restitution,
and contact/floor heights, re-anchored from the ball's real position at
every contact rather than integrated open-loop, so nothing drifts between
cycles. This matches the hybrid impact/flight model of
\cite{buehler1994,batz2009}; the first working version (09-15) had to fix
an early bug where the ball spawned at rest and never had enough energy to
complete its first bounce.

\subsection{Mobile Base Integration (Objective 2)}
The robot is a kinematic chassis connected to the paddle by a fixed-length
arm whose \emph{angle} is a free variable chosen independently each step,
rather than a rigid heading-plus-90$^\circ$ offset. The paddle is always
set to the ball's real current position (since they must physically
coincide for a contact to register), and the chassis trails it at arm's
length and the commanded angle, rendered as a debug line so the arm's pose
is visible during a run. This separation -- angle as a free variable --
turned out to be the key degree of freedom Objective~3's planner needed.

\subsection{Static Obstacle Course and Collision-Free Planning (Objective 3)}
This objective is complete and now exceeds its original scope: the
proposal specified 4--6 static obstacles, and the implementation generates
6--9 randomly placed and oriented crates, pillars, and walls per run. A
two-stage planner, run before the simulation starts, guarantees the whole
course is traversable with zero live obstacle checks during the run:
\begin{enumerate}[nosep]
  \item \textbf{Robot-body path search:} a hybrid-A* search over a state
  lattice of $(x, y, \text{heading})$ -- heading is part of the search
  state, not derived afterward, so the planner can angle the chassis
  through a gap that is only passable at an angle. Each edge is a short
  forward arc; the state being searched is the robot \emph{body's own
  center}, since the chassis is the part with the larger footprint that
  must clear every obstacle.
  \item \textbf{Arm-angle sequence search:} a second, Dijkstra-style search
  over the already-found path that picks a collision-free arm angle (ball
  position relative to the chassis) at every point and verifies that every
  \emph{transition} between consecutive points is collision-free, not just
  each point in isolation.
\end{enumerate}
If a randomly generated course turns out unsolvable at either stage,
generation retries with a fresh course (up to 50 attempts) rather than
failing outright. The two searches' near-identical priority-queue bookkeeping 
was unified into one shared search routine -- verified, across a ten-seed 
regression sweep, to produce byte-identical courses, paths, and arm-angle 
sequences to the pre-refactor implementation.

\section{Adaptations and Challenges}
Two design pivots came directly from discovering that an initially
reasonable approach could not give the guarantee the project needs.

\textbf{Reactive $\rightarrow$ proven-in-advance arm planning.} The first
version of the arm-angle controller chose the chassis's offset angle
live, frame by frame, reacting to obstacles as they were encountered, with
a durable-candidate lookahead and a wide safety margin around each
candidate angle. This could not actually \emph{guarantee} zero emergency
corrections: it only ever verified that some angle was clear \emph{at a
point}, never that the chassis could actually get from one point's safe
angle to the next's within the distance available. Three or more obstacles
narrowing the safe zone from multiple sides at once could still isolate
the tracked angle on an island with no reachable escape, however far ahead
the lookahead searched. We replaced it with the second search stage
described above, which plans the whole angle sequence up front and proves
every transition collision-free before the run starts -- the same
philosophy already used for the robot-body path itself.

\textbf{Ball-centric $\rightarrow$ robot-body-centric planning target.}
The first working planner searched the \emph{ball's} path and derived the
chassis's position as an offset from it. This was re-architected this week
so the search instead plans the robot \emph{body's} own center directly,
with the ball/paddle's target path derived from it afterward by solving
the same arm-offset relationship in reverse. The chassis is the part that
actually has to clear every obstacle by its full footprint, so planning
around it directly -- rather than planning around a smaller body and hoping
the larger one fits -- is a better match between what is being searched
and what the binding collision constraint actually is.

Separately, once the two-stage pipeline proved reliable, we increased
obstacle density beyond the proposal's 4--6 range to 6--9 to stress-test
the planner and produce a more demanding demonstration; Table~\ref{tab:prelim}
shows it still solves reliably at this density.

\section{Preliminary Results}
Table~\ref{tab:prelim} summarizes ten independently-seeded random courses
generated with the current 6--9-obstacle configuration, produced by
\texttt{planner\_benchmark.py} (\texttt{python planner\_benchmark.py
--latex} regenerates this exact table; \texttt{--n N} samples N fresh
random seeds instead). This is a sanity-check sweep run during
development, not the systematic density-vs-performance evaluation planned
for Objective~5. Course generation and both planning stages are pure
Python with no physics-engine dependency, so these numbers reflect planner
performance in isolation.

\begin{table}[h]
\centering
\small
\setlength{\tabcolsep}{4pt}
\begin{tabular}{@{}rrrrr@{}}
\toprule
\textbf{Seed} & \textbf{Obst.} & \textbf{Attempts} & \textbf{Path (m)} & \textbf{Time (s)} \\
\midrule
1       & 7 & 1 & 6.30 & 0.49 \\
2       & 6 & 1 & 6.30 & 0.38 \\
3       & 7 & 4 & 6.20 & 1.41 \\
42      & 6 & 3 & 6.30 & 1.04 \\
1000    & 8 & 4 & 6.30 & 1.15 \\
999999  & 7 & 1 & 6.30 & 0.18 \\
7       & 8 & 6 & 6.70 & 1.84 \\
55      & 8 & 2 & 6.30 & 0.33 \\
2024    & 6 & 8 & 6.20 & 2.24 \\
31415   & 8 & 1 & 6.30 & 0.19 \\
\bottomrule
\end{tabular}
\caption{Ten random seeds, 6--9 obstacles each, against a 6.00\,m
straight-line course. ``Attempts'' is regeneration attempts before a
solvable course was found (max 50 allowed); ``Time'' is total wall-clock
for course generation plus both planning stages.}
\label{tab:prelim}
\end{table}

All ten seeds solved well within the 50-attempt budget (mean 3.1 attempts,
worst case 8), with planned-path overhead over the straight-line course
ranging 3.3--11.7\% (mean 5.3\%) and total planning time under 2.3\,s per
course even in the worst case observed. This gives an early, favorable
signal that the two-stage guarantee is cheap enough in practice to keep
running at this obstacle density, not just sound in principle.

We have also added an automated MP4 recording mode (\texttt{--video}) that
captures a full run end-to-end for the demonstration component of this and
future deliverables.

\section{Schedule and Scope Update}
Table~\ref{tab:schedule} compares the original plan to actual progress.
We remain aligned with the checkpoint structure from the proposal; no
change to the Week~16 final deadline or the Week~14 buffer is proposed.

\begin{table}[h]
\centering
\small
\setlength{\tabcolsep}{4pt}
\begin{tabular}{@{}>{\raggedright\arraybackslash}p{1.05cm}p{6.35cm}@{}}
\toprule
\textbf{Wk} & \textbf{Status} \\
\midrule
3     & \textit{Done.} Proposal submitted; environment stood up. \\
4--5  & \textit{Done.} Working stationary dribble (Obj.~1), fixed an
energy bug in the first version. \\
6--7  & \textit{Done.} Mobile base dribbles while turning (Obj.~2);
static obstacle course and two-stage collision-free planner (Obj.~3),
later re-architected around the robot body and increased to 6--9
obstacles. \\
8     & \textbf{This submission.} Baseline dribble + static-obstacle
navigation demonstrated; code modularized. \\
9--10 & \textit{Planned, as proposed.} Redirect-the-bounce planner
(Obj.~4); see Remaining Work. \\
11--12 & \textit{Planned, as proposed.} Robustness passes; metrics/
logging pipeline (Obj.~5). \texttt{ContactEvent} logging is already in
place as an extension point. \\
13    & \textit{Planned, as proposed.} Full evaluation sweep across
obstacle densities. \\
14    & Buffer week, contingency unchanged from proposal. \\
15--16 & Finalize results, demo video, final report/presentation. \\
\bottomrule
\end{tabular}
\caption{Schedule status as of Week~8. No dates have moved.}
\label{tab:schedule}
\end{table}

\section{Remaining Work}
\textbf{Objective 4 -- Redirect-the-bounce planner.} The current system
already lets the \emph{robot} weave around obstacles and lets the
ball/paddle's offset angle vary to clear them, but it does not yet, as the
proposal specifies, shift the ball's own landing spot in response to a
blocked nominal target. Given the Section~III pivot toward proving
trajectories in advance rather than reacting live, our plan is to extend
the existing offline search -- rather than add a separate reactive
layer -- so that candidate landing spots become part of the same
collision-free search already proving the arm-angle sequence, preserving
the zero-live-correction guarantee instead of trading it away for a
reactive patch.

\textbf{Objective 5 -- Evaluation.} With Objective~4 in place, we will run
the density sweep proposed originally (dribble-survival rate, path-length
overhead, and redirect-event frequency vs.\ an obstacle-free baseline),
using the course-generation diagnostics already logged by the current
pipeline (obstacle count, attempts, path length, heading swing) as a
starting point for the metrics/logging pipeline.

\section*{Appendix: Individual Contributions}
\textit{Draft below, grounded in commit history; each author should review
and amend their own subsection before submission.}

\subsection*{Rudy Arsenault}
Created the initial environment to get the simulation working and implemented the first working
dribble test; diagnosed and fixed the energy bug that kept the ball's
first bounce from reaching hand height, producing the first reliable
stationary dribble (Objective~1); identified that the initial contact
model produced juggling-style behavior rather than a basketball-style
dribble and corrected the geometry to a two-event hybrid model --- an
actively controlled hand contact (paddle) above the ball and a passive,
restitution-based floor bounce below it; diagnosed and fixed a
corresponding object-placement overlap between the ball and paddle;
verified model produced a stable periodic dribble by logging the
ball's apex height and period across bounce cycles while testing.

\subsection*{Tracy Yang}
Built the mobile base integration (chassis + arm + paddle turning and
translating while dribbling, Objective~2); designed and implemented the
static obstacle course generator and the two-stage hybrid-A*/arm-angle
planner (Objective~3), including the pivot from reactive to proven-in-advance
arm planning and the later re-architecture to plan around the robot body
directly; increased obstacle density from the proposed 4--6 to 6--9;
modularized the codebase into separate concern-based files and removed
duplication between the two planning searches; added automated MP4 run
recording for demonstrations.

\begin{thebibliography}{9}
\bibitem{buehler1994} M.~Bühler, D.~E. Koditschek, and P.~J. Kindlmann,
``Planning and control of robotic juggling and catching tasks,''
\emph{The International Journal of Robotics Research}, vol.~13, no.~2,
pp.~101--118, 1994.

\bibitem{batz2009} G.~Bätz, M.~Sobotka, D.~Wollherr, and M.~Buss, ``Robot
basketball: Ball dribbling -- a modified juggling task,'' in \emph{Advances
in Robotics Research}, T.~Kröger and F.~M. Wahl, Eds. Berlin, Heidelberg:
Springer, 2009, pp.~323--334.
\end{thebibliography}

\end{document}
